\chapter{三相记忆的信息相态}
\chaptervoice{气、液、晶不是三种物理存储介质，也不是三个功能模块，而是描述认知结构变化速度、关系稳定度和可重构性的类比。$h$ 通常快速而松散，$E$ 保留流动轨迹，$\Theta$ 相对稳定并具有生成约束；三者仍可由同一种计算载体实现。}

\section{信息相态的操作性坐标}

为了避免物态隐喻滑向物理化，先给每个认知结构片段 $c$ 定义相态坐标
\begin{equation}
 \phi(c,t)=
 \bigl(\nu_c,\kappa_c,\tau_c,\rho_c\bigr),
 \label{eq:ch24-phase-coordinate}
\end{equation}
$\nu_c$ 是更新率，$\kappa_c$ 是关系稳定度，$\tau_c$ 是保持时间，$\rho_c$ 是在当前表示中的结构约束密度。

“气态、液态、晶态”是相态空间中的三个典型区域，而不是互斥标签。结构可能位于过渡区，也可能在某些坐标上稳定、另一些坐标上快速变化。分类阈值依任务、主体和观察窗口而定。

相态描述结构怎样存在和变化；第14章功能图描述它承担什么作用。二者是正交坐标，后文不能用“晶态”自动推断“长期记忆功能”，也不能用“当前规划功能”自动推断“气态”。

\section{$h$ 的气态类比}

$h$ 中结构受当前输入、目标和竞争持续重排，通常具有高更新率、短保持期和较低关系稳定度：
\begin{equation}
 \nu_h\ \text{高},\qquad
 \tau_h\ \text{短},\qquad
 \kappa_h\ \text{相对低}.
 \label{eq:ch24-gas}
\end{equation}
“气态”强调可快速组合、扩散和消散，不表示没有结构。

当前工作区仍有目标、权限和语法约束，否则无法形成连贯行动。一个高度确定的中间计算也可短暂存在于 $h$；它之所以属于 $h$，是因为仍处于当前可直接控制的状态，而非因为内容模糊。

气态结构的主要风险是干扰和错误绑定。新输入过多会挤出必要上下文，未完成关系可能在写入 $E$ 前消失。因此 $h$ 需要容量门、调度和事件边界。

\section{$E$ 的液态类比}

$E$ 将瞬时状态连接成带方向、时序和结果的轨迹。其关系比 $h$ 稳定，却仍能随新事件追加、分段、链接和重新解释：
\begin{equation}
 \nu_h>\nu_E>\nu_\Theta,\qquad
 \tau_h<\tau_E<\tau_\Theta
 \label{eq:ch24-rate-order}
\end{equation}
是典型而非绝对的时间尺度关系。

“液态”强调轨迹能够流动、聚合和重组。事件原件应保持，解释层则可随 $\Theta$ 更新。若把 $E$ 完全冻结，会失去情境重构；若允许随意改写原始观测，则会失去证据性。

$E$ 还承担相变缓冲：它阻止一次 $h$ 波动直接固化为 $\Theta$，又为生成结构提供跨样本证据。没有这一中间相，系统容易把短期噪声训练成长期偏差。

\section{$\Theta$ 的晶态类比}

$\Theta$ 中的关系经过多次经验和验证后相对稳定，具有更长保持时间、更低更新率和较高约束密度：
\begin{equation}
 \tau_h\ll\tau_E\ll\tau_\Theta,
 \qquad
 \kappa_h<\kappa_E<\kappa_\Theta.
 \label{eq:ch24-crystal}
\end{equation}
这些不等式表达典型尺度分离，而非所有元素逐项严格排序。

“晶态”强调稳定关系能约束并生成大量当前状态与轨迹。概念、技能和世界模型可用较少结构产生多种预测，这不同于保存大量事件副本。

晶态仍可发生缺陷、局部重排和相变。若环境改变而 $\Theta$ 完全冻结，稳定性会转为僵化；若更新过快，又会失去跨时生成能力。慢更新必须与现实残差和来源谱系连接。

\section{变化速度、稳定度与信息量不同}

更新慢不等于信息量大，稳定也不等于真实。一条错误规则可以长期稳定，一段短暂观测也可能含有关键新信息。相态轴不应用来评价真值、重要性或复杂度。

可分别定义变化速度
\begin{equation}
 \nu_c(W)=\frac{1}{|W|}
 \int_W d(c_{t+\Delta},c_t)\,dt
 \label{eq:ch24-change-rate}
\end{equation}
与稳定度
\begin{equation}
 \kappa_c(W)
 =1-\frac{\operatorname{Var}_W[OpSem(c_t)]}
          {Z_\kappa}.
 \label{eq:ch24-stability}
\end{equation}
前者测结构变化，后者测操作语义在窗口内的一致性。

若只按存储时长分类，缓存中的陈旧副本会被误称晶态；若只按访问频率分类，经常调用的稳定技能又会被误称气态。必须联合考察更新、保持和关系重构。

\section{相变是受门控的表示与动力学变化}

凝聚 $h\to E$、结晶 $E\to\Theta$、召回 $E\to h$、生成 $\Theta\to E/h$ 都是相态转换，但不要求数据在物理位置间移动。改变的可能是索引、关系密度、更新策略和可操作方式。

相变事件可写为
\begin{equation}
 Phase(c,t^+)
 =T_{pq}\bigl(Phase(c,t^-),Evidence,
              Goal,Authority\bigr),
 \label{eq:ch24-transition}
\end{equation}
$p,q$ 是源与目标区域。证据、目标和权限共同控制转换。

相变存在迟滞。刚形成的结构不应因一次反例立即完全熔解，旧结构也不能因历史稳定而拒绝持续反证。不同风险类别需要不同进入和退出门槛。

\section{相态轴与功能轴的二维表示}

令功能位置为 $f\in V_F$，信息相态为 $p\in\{h,E,\Theta\}$，则 Agent-CSO 的承载可写成二维分布
\begin{equation}
 \rho_C(f,p,t).
 \label{eq:ch24-two-axis}
\end{equation}
同一功能节点可同时处理三种相态，同一相态也可分布于多个功能节点。

例如规划功能可以读取 $\Theta$ 规则、召回 $E$ 轨迹并在 $h$ 组合当前方案；身体控制功能也可拥有快速活动状态、技能经历和稳定生成结构。因此三相不等于高层、中层、低层功能。

二维表达防止架构误配：相态转换由记忆动力学管理，功能迁移由责任重分配管理。一次技能学习可能同时发生 $E\to\Theta$ 结晶和规划到局部控制的 $M_F$ 迁移，但两者必须分别记录。

\section*{本章结论与进入下一章的接口}

气、液、晶只是在更新率、关系稳定度、保持时间和约束密度上的信息相态类比。典型关系是 $\tau_h\ll\tau_E\ll\tau_\Theta$，但相态不决定真值、重要性、物理载体或功能位置。三相转换受证据、目标、权限和迟滞门控，功能轴与相态轴必须二维表示。下一章将把三相放入生命周期结构时间，区分过去轨迹、现在切面和未来生成可能性。

\begin{readercheck}
分别找一个长期稳定的错误规则、短暂但关键的观测和高频调用的稳定技能；说明为何保持时间、真值、重要性与调用频率不能混为一个相态指标，并把三者放入功能—相态二维坐标。
\end{readercheck}

\cleardoublepage
